\section{局限、结论与建议}
\label{sec:limits}

\subsection{本次审计没有做到的事}

\begin{enumerate}
\item \textbf{没有取得 Nielsen 论文的全文。}PhilArchive／PhilPapers 对自动访问返回反爬验证页面（HTTP 403），我们没有绕过它；查询时互联网档案馆也没有存档。
本报告核查的主张来自：该文摘要被多处转述的句子、作者本人在 X 上的表述、以及 DEV Community 复核文章所引用的内容。
如果论文正文（据报道已更新十余版）包含未被这些来源反映的其他论据，本报告没有评估它们。本报告的结论针对的是公开流传的那条论证——即委托人截图中的那条。
\item \textbf{没有从源码重新编译整个证明库。}17.5 万个模块的完整构建在单机上无法在 7 小时预算内完成（公开的验证使用了云端多机构建）。
「导出文件确由这份源码生成」这一环我们依赖 MIT 页面所述的沙箱重建（重建所得导出与公布文件逐字节相同）；我们自己核对的是源码包与导出文件各自的哈希，以及挑战导出的独立重建。
\item \textbf{没有审读数学内容。}我们没有、也不可能在此时间内读懂 1740 万行证明的数学。本报告对「证明的数学思想是否自然、是否可以简化、几何在其中扮演什么角色」不作判断。
\item \textbf{可信基。}结论依赖于：Lean 的类型论是一致的；多个独立内核实现（官方内核、nanoda、lean4lean、con-leche、con-ron）不会同时在同一处出错；
\path{lean4export} 忠实地导出了环境；Mathlib 对 \lean{riemannZeta} 与 \lean{RiemannHypothesis} 的定义表达了通常意义的黎曼猜想。
其中最后一项是数学共同体的判断，不是本报告能机械验证的；MIT 页面提到的 Robin 不等式与 Lagarias 不等式推论提供了一个不依赖 $\zeta$ 定义的旁证，我们未复核。
\item \textbf{工具本身。}conescan 是为本次审计临时编写的（源码约 320 行，已提交供复核），其正确性依赖于「导出按依赖顺序排列」这一假设和若干自检；
con-leche 是一个很新的工具，其实现由 AI 编写并附形式化证明。我们把二者与第三方已公布的 nanoda、lean4lean 结果互相印证，而不单独依赖其中任何一个。
\item \textbf{交叉验证的独立性有限。}四个 \texttt{agy} 代理与主审计方阅读的是同一份源码；被本报告采用的代理结论均经复核，未复核者已在文中标明。
\end{enumerate}

\subsection{结论}

\begin{verdictbox}
\textbf{Nielsen 的打假不成立。}

\begin{itemize}
\item 其\textbf{事实观察}大体准确，本报告还把其中最关键的一条（\lean{anyGeometry} 与退化模型确在主定理依赖链上）从「未确认」提升为「内核级确认」，并纠正了一处细节（\lean{plugAll} 不在链上，在链上的是其孪生定理）。
\item 其\textbf{推理}不成立：对一个不带前提、类型为 Mathlib \lean{RiemannHypothesis}、只依赖三条标准公理、并经多个独立内核检查的闭定理而言，
「中间见证平凡」既不构成也不蕴含「证明无效」。在这份代码里，被指为「造假」的几何槽是一个只写不读的死参数。
\item 我们按四条可能的造假途径逐一排查，\textbf{没有找到这份证明的任何 bug 或 reward hacking}。就委托人「找出证明的 bug 就算赢」的目标而言，本次审计的如实答案是：没有找到，而且有充分理由相信在被审计的范围内不存在。
\end{itemize}
\end{verdictbox}

\subsection{建议}

\begin{itemize}
\item \textbf{给读者：评估「Lean 证明造假」类说法的五问清单。}
(1) 最终定理的陈述是不是你以为的那个命题（最好由只导入标准库的独立文件陈述并机械比对）？
(2) 它有没有前提？
(3) \texttt{\#print axioms}／导出文件里的公理是不是只有标准三条？
(4) 证明是否被独立于构建环境的内核重放过？
(5) 被检查的文件与发布的源码是否哈希对应？
五问都过关时，「中间定义看起来很可疑」不再是有效性问题；任何一问不过关，才值得深究。
\item \textbf{给证明作者与维护社区：}删除或显式注释这个未被使用的几何槽（\lean{Parameters.Γ}、\lean{OAt.geo}、\lean{GeoCl}）以及不在依赖锥内的「真几何」代码，
并公布一份「依赖锥内声明清单」。把 1,468,418 个真正被用到的声明与环境中其余约 123 万个常量区分开，是让人类开始理解这份证明的第一步，也能从源头消除此类误读。
\item \textbf{给批评者：}这个观察本身是有价值的代码考古——它准确地指出了形式化证明与其命名所暗示的数学叙事之间的落差。
把它表述为「该证明的几何部分是空壳，实际论证不经过它」是站得住的；把它表述为「该证明无效」则站不住。
\end{itemize}
